% GENERATED FILE. Edit canonical lessons, then run npm run generate:books.
% canonical-source-hash: fnv1a64:c5c1d4f4cdf1baf0
% canonical-lessons: GE-C273-veroeffentlichen, GE-C273-erwaehnen, GE-C273-aeussern, GE-C273-ankuendigen, GE-C273-bekanntgeben

\chapter{Bekannt machen}
\label{ch:ge-bekannt-machen}
\hlchaptermodality{273}

\begin{chapteropening}
\textbf{By the end of this chapter:} \emph{I can say how news was made known: to publish, to mention, to express, to announce and to make public.}

Complete the last lesson of chapter 273: I can say how news was made known: to publish, to mention, to express, to announce and to make public.
\end{chapteropening}

\section[veröffentlichen]{veröffentlichen — to publish}

\label{lesson:GE-C273-veroeffentlichen}



\begin{quote}
Before the new one: say the German for the journalist, then the German for the front page.
\end{quote}

\subsection*{You'll want to know: veröffentlichen}
\textbf{veröffentlichen} — \textquotedblleft{}to publish\textquotedblright{}.

\begin{quote}
Die Zeitung veröffentlichte das Foto am Montag.
\end{quote}

\begin{quote}
\emph{The newspaper published the photo on Monday.}
\end{quote}

Built on \emph{öffentlich}, already known. Perfekt: \emph{hat veröffentlicht}.

\begin{cousinweb}
\emph{veröffentlichen} is built on \emph{öffentlich} \textquotedblleft{}public\textquotedblright{}, from \emph{offen} \textquotedblleft{}open\textquotedblright{}, an inherited Germanic word; EtymDB lists no parent for the verb: to make open to all.
\end{cousinweb}

\begin{grammarlens}[title={What you've built}]
One more word for telling, explaining and describing at B1.
\end{grammarlens}

\subsection*{Your turn}
\begin{itemize}
\raggedright
  \item \emph{Say it:} \emph{veröffentlichen}
  \item \emph{Say it:} \emph{veröffentlichen}, once more
  \item \emph{Say it:} a sentence of your own with \emph{veröffentlichen}
  \item \emph{Recall:} say the German for the journalist, then the German for the front page, then say \emph{veröffentlichen} again
\end{itemize}

\begin{tcolorbox}[breakable,colback=teal!4,colframe=teal!35!black,title={Before you move on}]
What does veröffentlichen mean? (\textquotedblleft{}To publish\textquotedblright{}.) Say it once more.
\end{tcolorbox}
\section[erwähnen]{erwähnen — to mention}

\label{lesson:GE-C273-erwaehnen}



\begin{quote}
Before the new one: say the German for the front page, then the German for to publish.
\end{quote}

\subsection*{You'll want to know: erwähnen}
\textbf{erwähnen} — \textquotedblleft{}to mention\textquotedblright{}.

\begin{quote}
Er erwähnte den Unfall mit keinem Wort.
\end{quote}

\begin{quote}
\emph{He did not say a single word about the accident.}
\end{quote}

It takes a direct object: \emph{Sie erwähnte ihn nicht}.

\begin{cousinweb}
\emph{erwähnen} is built on \emph{wähnen} \textquotedblleft{}to suppose\textquotedblright{}, from Old High German \emph{wānen} (EtymDB), a cousin of old English \emph{ween} \textquotedblleft{}to think\textquotedblright{}.
\end{cousinweb}

\begin{grammarlens}[title={What you've built}]
One more word for telling, explaining and describing at B1.
\end{grammarlens}

\subsection*{Your turn}
\begin{itemize}
\raggedright
  \item \emph{Say it:} \emph{erwähnen}
  \item \emph{Say it:} \emph{erwähnen}, once more
  \item \emph{Say it:} a sentence of your own with \emph{erwähnen}
  \item \emph{Recall:} say the German for the front page, then the German for to publish, then say \emph{erwähnen} again
\end{itemize}

\begin{tcolorbox}[breakable,colback=teal!4,colframe=teal!35!black,title={Before you move on}]
What does erwähnen mean? (\textquotedblleft{}To mention\textquotedblright{}.) Say it once more.
\end{tcolorbox}
\section[äußern]{äußern — to express, to voice; reflexive: to comment}

\label{lesson:GE-C273-aeussern}



\begin{quote}
Before the new one: say the German for to publish, then the German for to mention.
\end{quote}

\subsection*{You'll want to know: äußern}
\textbf{äußern} — \textquotedblleft{}to express, to voice; reflexive: to comment\textquotedblright{}.

\begin{quote}
Sie äußerte ihre Meinung ganz offen.
\end{quote}

\begin{quote}
\emph{She voiced her opinion quite openly.}
\end{quote}

Reflexive, it means to comment: \emph{Er wollte sich nicht äußern}.

\begin{cousinweb}
\emph{äußern} comes from Middle High German \emph{iuzern} (EtymDB), built on \emph{außen} \textquotedblleft{}outside\textquotedblright{}, Old High German \emph{ūzana}: to bring something out.
\end{cousinweb}

\begin{grammarlens}[title={What you've built}]
One more word for telling, explaining and describing at B1.
\end{grammarlens}

\subsection*{Your turn}
\begin{itemize}
\raggedright
  \item \emph{Say it:} \emph{äußern}
  \item \emph{Say it:} \emph{äußern}, once more
  \item \emph{Say it:} a sentence of your own with \emph{äußern}
  \item \emph{Recall:} say the German for to publish, then the German for to mention, then say \emph{äußern} again
\end{itemize}

\begin{tcolorbox}[breakable,colback=teal!4,colframe=teal!35!black,title={Before you move on}]
What does äußern mean? (\textquotedblleft{}To express, to voice; reflexive: to comment\textquotedblright{}.) Say it once more.
\end{tcolorbox}
\section[ankündigen]{ankündigen — to announce}

\label{lesson:GE-C273-ankuendigen}



\begin{quote}
Before the new one: say the German for to mention, then the German for to express.
\end{quote}

\subsection*{You'll want to know: ankündigen}
\textbf{ankündigen} — \textquotedblleft{}to announce\textquotedblright{}.

\begin{quote}
Die Firma kündigte neue Preise an.
\end{quote}

\begin{quote}
\emph{The company announced new prices.}
\end{quote}

Separable: \emph{sie kündigt an}, \emph{sie hat angekündigt}.

\begin{cousinweb}
\emph{ankündigen} is built on \emph{kündigen}, from \emph{kund} \textquotedblleft{}known\textquotedblright{}, an inherited word with the Gothic cousin \emph{kunþs} \textquotedblleft{}known\textquotedblright{} (EtymDB): to make known ahead.
\end{cousinweb}

\begin{grammarlens}[title={What you've built}]
One more word for telling, explaining and describing at B1.
\end{grammarlens}

\subsection*{Your turn}
\begin{itemize}
\raggedright
  \item \emph{Say it:} \emph{ankündigen}
  \item \emph{Say it:} \emph{ankündigen}, once more
  \item \emph{Say it:} a sentence of your own with \emph{ankündigen}
  \item \emph{Recall:} say the German for to mention, then the German for to express, then say \emph{ankündigen} again
\end{itemize}

\begin{tcolorbox}[breakable,colback=teal!4,colframe=teal!35!black,title={Before you move on}]
What does ankündigen mean? (\textquotedblleft{}To announce\textquotedblright{}.) Say it once more.
\end{tcolorbox}
\section[bekanntgeben]{bekanntgeben — to make public, to announce officially}

\label{lesson:GE-C273-bekanntgeben}



\begin{quote}
Before the new one: say the German for to express, then the German for to announce.
\end{quote}

\subsection*{You'll want to know: bekanntgeben}
\textbf{bekanntgeben} — \textquotedblleft{}to make public, to announce officially\textquotedblright{}.

\begin{quote}
Die Schule gab das Ergebnis am Freitag bekannt.
\end{quote}

\begin{quote}
\emph{The school made the result public on Friday.}
\end{quote}

Separable: \emph{sie gibt bekannt}, \emph{sie hat bekanntgegeben}.

\begin{cousinweb}
\emph{bekanntgeben} joins \emph{bekannt} \textquotedblleft{}known\textquotedblright{} and \emph{geben} \textquotedblleft{}to give\textquotedblright{}, from Old High German \emph{geban}, Proto-Germanic \emph{gebaną} (EtymDB): to give out as known.
\end{cousinweb}

\begin{grammarlens}[title={What you've built}]
One more word for telling, explaining and describing at B1.
\end{grammarlens}

\subsection*{Your turn}
\begin{itemize}
\raggedright
  \item \emph{Say it:} \emph{bekanntgeben}
  \item \emph{Say it:} \emph{bekanntgeben}, once more
  \item \emph{Say it:} a sentence of your own with \emph{bekanntgeben}
  \item \emph{Recall:} say the German for to express, then the German for to announce, then say \emph{bekanntgeben} again
\end{itemize}

\begin{tcolorbox}[breakable,colback=teal!4,colframe=teal!35!black,title={Before you move on}]
What does bekanntgeben mean? (\textquotedblleft{}To make public, to announce officially\textquotedblright{}.) Say it once more.
\end{tcolorbox}
